\section{Exercise 1 - Stream Ciphers}

\begin{frame}{Linear Feedback Shift Register}
    In this exercise we are going to study the M-sequence generated by the LFSR presented in the following figure.

    \begin{figure}[h]
		\centering
		\includegraphics[width=.75\textwidth]{images/LFSR_schema.png}
	\end{figure}

    The LFSR is characterized by:

    \begin{itemize}
        \item $m = 10$ cells;
        \item polynomial description $D^{10} + D^{3} + 1$.
    \end{itemize}
\end{frame}

\begin{frame}{M-sequence}
    We choose an arbitrary seed ($0000010100$) to initialize the $m = 10$ cells of the register. Since the polynomial is primitive, the LFSR generates a periodic M-sequence of $N = 2^m - 1 = 1023$ bits. We denote this sequence as $v(n)$ and the corresponding bipolar sequence as $v'(n)$. In the following image we have plotted the first 20 bits of the bipolar sequence $v'(n)$ we obtained.
\end{frame}

\begin{frame}
    \begin{figure}[h]
		\centering
		\includegraphics[width=1\textwidth]{images/bipolar_seq.png}
	\end{figure}
\end{frame}

\begin{frame}{Comments}
    From the figure you can see that the first 10 bits are simply the bits of the seed in reverse order, because the output bit is always the rightmost bit of the register state.
    
    Furthermore, the M-sequence $v(n)$ is independent of the seed we choose, in fact a different initialization (except for the all-zero seed) would have produced only a shifted version of $v(n)$. The M-sequence depends only on the polynomial description of the LFSR.
\end{frame}

\begin{frame}{Properties of the M-sequence - $0/1$}
    Now we will analyze a couple of properties of $v(n)$ to test the goodness of this pseudo-random sequence.

    We start by counting the number of bits that are $1$ or $0$, $N_1$ and $N_0$ respectively; the closer these two numbers are, the better the randomness. We get that $N_1 = 512$ and $N_0 = 511$, so the best case scenario. This is not a lucky coincidence, in fact we can prove that for any M-sequence the equation $N_1 = N_0 + 1$ is always verified.
\end{frame}

\begin{frame}{Demonstration - $0/1$}
    An M-sequence, or a sequence with maximum period, is generated by a Linear Feedback Shift Register using a primitive polynomial. The sequence has a period of $2^m - 1$, where $m$ is the number of cells in the LFSR. This period occurs because the shift register cycles through all possible states except the state where all bits are zero (in which case the sequence would remain static). Since the output bit, which will compose the sequence, is always the rightmost bit of the state, it will be $1$ in $2^{m-1}$ cases and it will be $0$ in $2^{m-1} - 1$ cases, because the all-zero state is excluded. So we have:
    \[
        N_0 = 2^{m-1} - 1
    \]
    \[
        N_1 = 2^{m-1} = N_0 + 1
    \]
    \[
        N_0 + N_1 = 2^{m-1} - 1 + 2^{m-1} = 2\cdot2^{m-1} - 1 = 2^{m} -1 = N
    \]
\end{frame}

\begin{frame}{Properties of the M-sequence - Runs}
    The next properties we will analyze are those concerning the runs. A run of length $i$ is any sequence of $i$ consecutive $1$s or $0$s in the $v(n)$. Let $\mathcal{N}_i$, $\mathcal{N}_i(0)$ and $\mathcal{N}_i(1)$ be the number of runs of length i, the number of runs of zeros of length i and the number of runs of ones of length $i$ in $v(n)$ respectively, we obtain the following table.
\end{frame}

\begin{frame}
    \begin{figure}[h]
		\centering
		\includegraphics[width=.54\textwidth]{images/runs_table.png}
	\end{figure}
\end{frame}

\begin{frame}{Properties of the M-sequence - Runs}
    By adding all the values of the last column of the table we can calculate the total number of runs, that is: $\mathcal{N}_T = 512$. Therefore, the values reported in the table verify all the run properties of an ideal binary random sequence:

    \begin{enumerate}
        \item $\mathcal{N}_T = \lceil \frac{N}{2} \rceil$, where $N$ is the length of the $v(n)$;
        \item $\mathcal{N}_i = \frac{\mathcal{N}_T}{2^i}$;
        \item $\mathcal{N}_i(0) = \mathcal{N}_i(1)$, except for $i = 9$ and $i = 10$.
    \end{enumerate}
\end{frame}

\begin{frame}{Demonstration - Runs}
    From the table above you can also see that there are no runs of zeros of length $m = 10$ and there are no runs at all of length greater than $m$. 
    
    We can prove that a run of zeros of length $m$ cannot be part of an M-sequence. Indeed, if $m$ consecutive zeros appeared, the LFSR would enter the all-zero state. Once in this state, the LFSR cannot leave it because the feedback is zero, which contradicts the fact that an M-sequence is periodic.

    We can also prove that an M-sequence cannot have a run of ones of length $m+1$. In fact, the feedback of the all-ones state must be $0$, otherwise, if the all-ones state produces a feedback of $1$, the LFSR would remain in that state indefinitely, breaking the periodicity of the sequence. Thus, after $m$ ones there will always be a zero.
\end{frame}

\begin{frame}{Properties of the M-sequence - Autocorrelation}
    The last property of the M-sequence we are going to analyze is the autocorrelation function:
    \[
        R(\tau) = \sum_{n=0}^{N-1}v'(n)v'(n-\tau)
    \]
    with $-\frac{N-1}{2} \le \tau \le \frac{N-1}{2}$.

    As you can see in the following graph, the function verifies the properties of an ideal random sequence:

    \begin{itemize}
        \item $R(\tau) = N$, for $\tau = 0$
        \item $R(\tau) = -1$, for $\tau \ne 0$
    \end{itemize}
\end{frame}

\begin{frame}
    \begin{figure}[h]
		\centering
		\includegraphics[width=.55\textwidth]{images/autocorrelation.png}
	\end{figure}
\end{frame}

\begin{frame}{Breaking a Linear LFSR}
    By a chosen plaintext attack, we manage to intercept this portion of the key: $1010100010000001101010010111011000101110$ (40 bits).
    We now want to compute the primitive polynomial of the LFSR that generated the key.

    To do so, we use the known portion of the key to compute the polynomial description of the LFSR for an increasing number of cells, starting at $m = 3$ (the maximum number of cells is 20). For each polynomial found, we calculate a 40-bit sequence and compare it with the portion of the key we have. If the two sequences are equal and the polynomial is primitive, then we have found the correct polynomial description.
\end{frame}

\begin{frame}{Breaking a Linear LFSR}
    The primitive polynomial found by the code is as follows:
    \[
    p(D) = D^{19} + D^5 + D^2 + D + 1
    \]
    and the number of cells in the LFSR is $m = 19$.
    
    We can now compute the next 20 bits of the key, which are: $00010000111000010000$.
\end{frame}

\begin{frame}{Nonlinear-Feedback Shift Register}
    The main weakness of LFSRs is that the feedback function is linear. By exploiting this linearity, an attacker can decrypt the ciphertext with only a small portion of the key ($2m$ bits is enough). To overcome this problem, modern stream encryption algorithms use nonlinear feedback shift registers (NLFSRs), which have nonlinear feedback functions.

    NLFSRs are much more complex than LFSRs, so while a general method for generating maximal-length sequences is known for LFSRs (e.g. using primitive polynomials), the same is not true for NLFSRs. However, some feedback functions for NLFSRs are known that are able to produce long periods, at least for small values of the number of cells ($m \le 25$).
\end{frame}

\begin{frame}{Starting LFSR}
    The following schema represents the LFSR that we want to improve by transforming it into an NLFSR. His polynomial description is $p(D) = D^5 + D + 1$, which is a primitive polynomial.

    \begin{figure}[h]
		\centering
		\includegraphics[width=.7\textwidth]{images/lfsr.png}
	\end{figure}    
\end{frame}

\begin{frame}{Showing linearity}
    In the case of the LFSR, given 2 different seeds:
    \[
    S1 = 11111 \quad S2 = 10101
    \]
    and their corresponding M-sequences:
    \[
    M1 = 1111100110100100001010111011000
    \]
    \[
    M2 = 1010111011000111110011010010000
    \]
    we have that the M-sequence generated by the sum of the seeds
    \[
    S3 = S1 \oplus S2 = 01010
    \]
    is equal to the sum of the M-sequences
    \[
    M3 = M1 \oplus M2 = 0101011101100011111001101001000
    \]
\end{frame}

\begin{frame}{Adding nonlinearity}
    A very common solution to add nonlinearity in the feedback function is to insert logic gates such as AND and OR in the schematic. In our case, an AND gate will be inserted between the output of cell 1 and the output of cell 3 and the result will be added to the feedback along with the value of cell 1.

    \begin{figure}[h]
		\centering
		\includegraphics[width=.7\textwidth]{images/nlfsr.png}
	\end{figure}  
\end{frame}

\begin{frame}{Resulting NLFSR}
    The obtained NLFSR has the peculiarity of having a period of $2^m-1 = 31$ bits, that is, the maximum period. The feedback function can be written as:

    \[
    f(x_0, x_1, x_2, x_3, x_4) = x_0 \oplus x_1 \oplus x_3 \oplus x_1 x_3
    \]
\end{frame}

\begin{frame}{Showing nonlinearity}
    In the case of the NLFSR, given 2 different seeds:
    \[
    S1 = 11111 \quad S2 = 10101
    \]
    and their corresponding M-sequences:
    \[
    M1 = 1111100110100100001010111011000
    \]
    \[
    M2 = 1010111011000111110011010010000
    \]
    we have that the M-sequence generated by the sum of the seeds
    \[
    S3 = S1 \oplus S2 = 01010
    \]
    \[
    M3 = 0101011101001101100100011111000
    \]
    is different from the sum of the M-sequences:
    \[
    M4 = M1 \oplus M2 = 0101011011001100011011100111000
    \] 
\end{frame}

\begin{frame}{References}
    Source for NLFSR: \href{https://www.diva-portal.org/smash/get/diva2:812512/FULLTEXT01.pdf}{A List of Maximum Period NLFSRs - Elena Dubrova}.
\end{frame}